% ===== §10 =====
\section{The Descent of the Object at the Gap}
\label{sec:descent}

\noindent
This section establishes what takes form at the gap. The aim is to show that the objet petit a is not an object that fills the lack but the form the lack takes within the economy of desire, the shape manque assumes so as to set desire turning, and that it holds its place precisely by not being the filling of it. The section states the temptation to misread the object, gives the correct account with its psychoanalytic warrant, follows the object through its several forms as the partial objects around which the drive turns, reviews the reading that would treat it as an attainable thing, and draws the consequence that prepares the distinction of drive from desire.

\subsection*{What descends, and what it is not}

At the gap something takes form. It is not a thing that fills the gap, and not a content that closes it. It is the way the gap is given a shape within the economy of desire, and psychoanalysis calls it the objet petit a, the object-cause of desire \citep{lacan1977xi}. Everything in the account depends on holding the object apart from the ordinary notion of an object, and the temptation to collapse the two is strong, since the word object invites it.

\subsection*{Object-cause, not object of desire}

The objet a is not the object that desire seeks, the target it aims at and might attain. It is the object-\emph{cause}, the point from which desire is set in motion, and it stands in the place of manque. Lacan drew the distinction sharply in the seminar on the four fundamental concepts: the object around which the drive turns is not the object that would satisfy a need but the object-cause that sustains the turning, and the drive's circuit is defined by this object as an absence it circles rather than a presence it reaches \citep{lacan1977xi}. Žižek gives the point its sharpest formulation in commentary: the objet a is not a positive entity but the void of desire given a form, an object that is really only the materialization of a lack, such that what looks like the cause of desire is the lack itself wearing the mask of an object \citep{zizek1989sublime}. The gap itself is without form, a pure remainder, a want. The objet a is the form the want takes when it enters the field of desire, the shape manque assumes so as to set desire turning. It does not answer the lack. It gives the lack a figure, and by giving it a figure lets it function as the pivot of a movement.

\subsection*{The forms of the object}

The object-cause is not one thing but appears in several forms, and following them makes its character clearer, since what the forms have in common is exactly what distinguishes the object-cause from any ordinary object. Psychoanalysis lists among these forms not only the objects of the early drives but, in Lacan's extension, the gaze and the voice \citep{lacan1977xi}. What is striking is that these are not objects one grasps at all. The gaze is not what I see but the sense of being seen from a point I cannot occupy, a point in the field of vision that looks back and that I can never bring into my own view. The voice is not what is said, not the meaning carried, but a remainder of the vocal that no meaning absorbs, what persists in a voice when everything sayable in it has been understood. Each of these is an object that is constitutively out of reach, given only as something that escapes, and this is precisely why they can serve as forms of the object-cause. What the several forms share is not a property one could hold but a manner of escaping, and it is the escaping, not any content, that makes them able to figure the lack and set desire turning. The object-cause is whatever, in a field, functions as its unreachable point.

\begin{casebox}{the look one cannot return}
Between two people there is often a way the one feels seen by the other that they cannot reciprocate on demand, a sense of being held in a regard whose exact point they can never occupy or verify. It is not a fact one could confirm by asking, since to ask is already to fall out of it. This being-seen from a point one cannot reach is the gaze as object-cause, a form of the lack around which the bond turns, given only as what escapes the very look that would secure it.
\end{casebox}

\subsection*{Against the attainable object}

Read as an ordinary object, the objet a would be something one could reach and hold, and desire would be the project of attaining it. The whole clinical and theoretical weight of the concept stands against this reading, and the forms just surveyed show why. The gaze and the voice cannot be attained; they are given only as what escapes. Were the objet a in general attained and held, the gap it forms would be closed, and the gap is what gave the object its office. Attained, it would cease to be the object-cause and become a mere thing; and desire, finding it a mere thing, would pass on, because what desire turned around was never the thing but the want the thing had been figuring. This is the familiar structure of the object that disappoints in the having, whose possession reveals that it was never the possession that was wanted. The objet a keeps its power only so long as it remains the form of a lack and not the lack's cancellation.

\subsection*{The consequence}

Two things follow, and they organize the next section. First, because the object is the form of a lack and not an attainable thing, the movement that turns around it is a movement that does not aim to end, a circling rather than a reaching. Second, because no signifier and no thing coincides with the want the object figures, there is also a movement that runs along the signifiers without resting, a sliding rather than a circling. The one is drive, the other desire, and both turn on the object that descends at the gap without ever closing it. The next section separates them and shows that relational life needs both.
